\songtitle{\textjapanese{ビル・パイ ―雪と薔薇―} (Bill Pi: Snow \& Roses)}

\tag*{1989}\tag{punk-rock}\tag{japanese-lyrics}\tag{bill-pi}
% -> stories.tex (Named Fictional Works?)
\begin{notes}
Lyrics by Claude based on Carlos Thompson's 1989 story \emph{La historia de Bill Pi}, part one \emph{El regreso}, written as an anime-style theme song. Retrieved from Gabriel's Suno page (V6-MINI, 15 November 2026).
\end{notes}
\begin{style}
cinematic punk rock; taiko drums fire short, syncopated strikes beneath orchestral strings, distorted guitars and 90s retro synth stabs; driving half-time verses accelerate into double-time, huge build choruses; gritty male lead in a low, dark lower register delivers clipped, rhythmic staccato Japanese vocals, with a close-mic female whisper and layered gang vocals; rain and payphone dial tone give way to cymbal crashes and a hard stop.
\end{style}
\begin{lyrics}
\songpart{Intro - spoken, low whisper, rain, payphone dial tone}
\begin{japanese}
一九九〇年、ボゴタ。\\
二年ぶりに帰った男を、誰も見分けられない。
\end{japanese}

\songpart{Build - taiko hits, strings swell, guitar feedback}

\songpart{Verse 1 - gritty male vocal, half-spoken}
\begin{japanese}
モヒカンに 鎖と鋲\\
スイスの空は 嘘で凍ってた\\
ジェーンの笑顔は 毒入りの粉\\
五十三万ポンド 奪って逃げた
\end{japanese}

\songpart{Pre-Chorus - rising, drums double-time}
\begin{japanese}
パリ、ボゴタ、ニューヨーク\\
最後のコインで 呼び出す声
「誰だお前は?」友が囁く
\end{japanese}

\songpart{Chorus - full band, anthemic, gang vocals}
\begin{japanese}
ビル・パイ! ビル・パイ! 夜を駆け抜けろ\\
雪を売れ 薔薇を抱け\\
ビル・パイ! ビル・パイ! 振り向くな\\
追ってくるのは 料理人 (Le Cuisinier)
\end{japanese}

\songpart{Verse 2 - male vocal, tense, sparse drums}
\begin{japanese}
懐に隠す ナイフの冷たさ\\
赤いフィエスタ フレンチの夜\\
笑って 泣いて 触れて 離れて
「俺はもう あの頃の俺じゃない」
\end{japanese}

\songpart{Pre-Chorus - rising}
\begin{japanese}
ロンドン、ベルン、プラハ\\
受話器の向こう 銃声の影
\end{japanese}

\songpart{Chorus - key change up, full band, gang vocals}
\begin{japanese}
ビル・パイ! ビル・パイ! 夜を駆け抜けろ\\
雪を売れ 薔薇を抱け\\
ビル・パイ! ビル・パイ! 振り向くな\\
愛と裏切りの コード・ネーム
\end{japanese}

\songpart{Bridge - breakdown, female whisper, single piano}
\begin{japanese}
愛してる 切れる電話\\
金庫の鍵は テープの裏に\\
もし私が死んだなら——\\
その時 扉を開けて
\end{japanese}

\annotation{Spoken - trailer card, deep voice, heartbeat bass}
\begin{japanese}
「彼が死ねば、真実が開く」
\end{japanese}

\songpart{Final Chorus - maximum intensity, double vocals, cymbal crashes}
\begin{japanese}
ビル・パイ! ビル・パイ! 雪と薔薇の\\
最後の夜が 始まる
\end{japanese}

\songpart{Outro - spoken, strings fade, one guitar stab}
\begin{japanese}
ノート三冊。カセット一本。\\
ページをめくるたび、リリアナは——憎んでいく。
\end{japanese}

\annotation{Hard stop}
\annotation{Title card voice: whispered}\\
\begin{japanese}
ビル・パイ ―雪と薔薇―
\end{japanese}
\end{lyrics}

